%---
% title: "What is Delay Flow Matching doing?"
% date: 2026-07-10
% work_order: 4
% permalink: /posts/2026/07/delay-flow-matching-learns-the-coupling/
% tags:
%   - machine-learning
%   - flow-matching
%   - generative-models
%   - theory
%---
\documentclass[11pt]{article}
\usepackage{publication-web}
\begin{document}

Ordinary flow matching learns a marginal vector field from a sampled endpoint coupling. Delay Flow Matching (DFM) additionally conditions on an earlier point from the same constructed trajectory. For common deterministic interpolants, the two trajectory points identify the endpoint pair; DFM therefore learns a continuation conditioned on the sampled coupling, much like source-conditioned flow matching or Augmented Bridge Matching.

\section{The two objectives}\label{the-two-objectives}

Let the coupling pair source \(x_0\) with target \(x_1\), and construct

\[
x_t=\alpha_tx_0+\beta_tx_1,
\qquad
u_t=\dot\alpha_tx_0+\dot\beta_tx_1.
\]

Ordinary FM minimizes squared error for \(v_\theta(t,x_t)\), so its population solution is

\[
v^*_{\mathrm{FM}}(t,x)=\mathbb E[u_t\mid x_t=x].
\]

Different endpoint pairs crossing \(x\) are averaged. DFM instead learns \(v_\theta(t,x_t,x_{t-\tau})\), with

\[
v^*_{\mathrm{DFM}}(t,x,y)
=\mathbb E[u_t\mid x_t=x,\ x_{t-\tau}=y].
\]

The extra conditioning changes the estimand, not just the network input.

\section{Two points recover the endpoints}\label{two-points-recover-the-endpoints}

Write \(s=t-\tau\). If

\[
\Delta_{t,s}=\alpha_t\beta_s-\alpha_s\beta_t\neq0,
\]

then the coordinatewise \(2\times2\) system is invertible:

\[
x_0=\frac{\beta_sx_t-\beta_tx_s}{\Delta_{t,s}},
\qquad
x_1=\frac{-\alpha_sx_t+\alpha_tx_s}{\Delta_{t,s}}.
\]

Consequently, the DFM target is the deterministic continuation

\[
u_t(x_t,x_s)=
\frac{(\dot\alpha_t\beta_s-\dot\beta_t\alpha_s)x_t
+(\alpha_t\dot\beta_t-\beta_t\dot\alpha_t)x_s}
{\Delta_{t,s}}.
\]

For the straight schedule \(x_r=(1-r)x_0+rx_1\), this reduces to

\[
u_t=x_1-x_0=\frac{x_t-x_{t-\tau}}{\tau}
\qquad (t>\tau).
\]

The per-sample DFM gradient therefore regresses toward the continuation attached to one endpoint pair, whereas ordinary FM regresses toward \(\mathbb E[u_t\mid x_t]\). This is the same statistical distinction studied by conditional FM and coupling-preserving bridge objectives (\href{https://arxiv.org/abs/2210.02747}{Lipman et al.~2022}; \href{https://arxiv.org/abs/2302.00482}{Tong et al.~2023}; \href{https://arxiv.org/abs/2311.06978}{De Bortoli et al.~2023}).

\section{Two consequences}\label{two-consequences}

First, an ideal fitted model with an endpoint-determined coupling does not need a long numerical rollout to reveal its implied endpoint. At time zero, let

\[
\widehat u_0=v_\theta(0,x_0,x_0)
\approx\dot\alpha_0x_0+\dot\beta_0x_1.
\]

If \(\dot\beta_0\neq0\), then

\[
\widehat x_1=
\frac{\widehat u_0-\dot\alpha_0x_0}{\dot\beta_0}.
\]

For a straight schedule, \(\widehat x_1=x_0+\widehat u_0\); subsequent target vectors follow from the known schedule. Extra evaluations can still correct finite-capacity or numerical error, but they add no endpoint information in this idealized setting.

Second, the coupling becomes a modeling choice. With a product coupling, DFM preserves arbitrary noise--data pairings rather than the ordinary marginal FM field. A meaningful coupling---paired observations, optimal transport, minibatch matching, or a bridge construction---gives the extra input useful structure to preserve (\href{https://arxiv.org/abs/2304.14772}{Pooladian et al.~2023}; \href{https://arxiv.org/abs/2303.16852}{Shi et al.~2023}).

\section{Where the argument does not apply}\label{where-the-argument-does-not-apply}

In a genuine delay differential equation,

\[
\dot x(t)=f(x(t),x(t-\tau)),
\]

the earlier state is physical state, not a code for two interpolation endpoints. DFM may then supply variables required by the dynamics; the endpoint-inversion argument applies only when trajectories are constructed from endpoint pairs (\href{https://arxiv.org/abs/2102.10801}{Zhu et al.~2021}; \href{https://openreview.net/forum?id=6lH1XblLpo}{Zhao et al.~2026}).

\end{document}
